\documentclass[11pt,a4paper]{ctexart}

\usepackage[a4paper,margin=2.35cm]{geometry}
\usepackage{fontspec}
\usepackage{setspace}
\usepackage{indentfirst}
\usepackage{booktabs}
\usepackage{tabularx}
\usepackage{array}
\usepackage{enumitem}
\usepackage{xcolor}
\usepackage{hyperref}
\usepackage{titlesec}
\usepackage{fancyhdr}
\usepackage{lastpage}

\setmainfont{Times New Roman}
\setsansfont{Times New Roman}
\setmonofont{Consolas}
\setCJKmainfont[AutoFakeBold=2.5,AutoFakeSlant=0.18]{SimSun}
\setCJKsansfont[AutoFakeBold=1.8]{STFangsong}
\setCJKmonofont{SimSun}

\hypersetup{
  colorlinks=true,
  linkcolor=black,
  urlcolor=black,
  citecolor=black,
  pdftitle={MT5 + TSLib Quant Written Test},
  pdfauthor={Codex}
}

\setlength{\parindent}{2em}
\setlength{\parskip}{0.35em}
\setstretch{1.18}
\setlist[enumerate]{leftmargin=2.2em,itemsep=0.25em,topsep=0.25em}
\setlist[itemize]{leftmargin=2em,itemsep=0.25em,topsep=0.25em}

\definecolor{titlegray}{RGB}{60,60,60}

\titleformat{\section}
  {\Large\bfseries\color{titlegray}}
  {\thesection.}{0.6em}{}
\titleformat{\subsection}
  {\large\bfseries}
  {\thesubsection}{0.6em}{}

\pagestyle{fancy}
\fancyhf{}
\fancyhead[L]{\small MT5 + TSLib 研究答卷}
\fancyhead[R]{\small \thepage/\pageref*{LastPage}}
\renewcommand{\headrulewidth}{0.3pt}
\renewcommand{\footrulewidth}{0pt}

\newcolumntype{C}[1]{>{\centering\arraybackslash}p{#1}}

\begin{document}

\begin{center}
  {\zihao{2}\bfseries 基于 MT5 与清华 TSLib 的时间序列交易研究\par}
  \vspace{0.7em}
  {\small 数据源：MT5 模拟账户历史行情 \quad 模型库：Time-Series-Library (TSLib)\par}
  \vspace{0.25em}
  {\small 复现实验仓库：\url{https://github.com/GoodNight0721/Time-Series-Library}\par}
\end{center}

\vspace{0.6em}
\hrule
\vspace{0.9em}

\section{任务与修订方向}

题目要求将 MT5 历史行情与 TSLib 结合，训练时间序列模型并评估胜率、盈亏比和最大回撤。第一版实验已经完成了数据导出、模型训练和回测闭环，但基础策略采用的是“单根 K 线开平仓 + 低阈值触发”的最简实现，换手偏高，净收益偏弱。

根据反馈，第二版策略重点改了四个方向：
\begin{enumerate}
  \item 提高开仓阈值，减少低置信度交易；
  \item 改为状态化持仓，同方向预测连续出现时不重复换手；
  \item 增加趋势过滤与止损机制；
  \item 加入平仓后的冷却期，进一步压低无效交易频率。
\end{enumerate}

第二版的目标不再只是展示代码闭环，而是验证这些交易层修改能否把净收益从负值拉到正值。

\section{数据与模型}

数据来自 MT5 模拟账户，覆盖三类资产：
\begin{itemize}
  \item 外汇：\texttt{EURUSD}、\texttt{GBPUSD}、\texttt{USDJPY}
  \item 贵金属：\texttt{XAUUSD}、\texttt{XAGUSD}
  \item 指数 CFD：\texttt{US500}、\texttt{USTEC}
\end{itemize}

周期覆盖为外汇/黄金使用 \texttt{M15/H1/H4}，白银/指数使用 \texttt{H1/H4}。数据切分继续沿用 TSLib 的 \texttt{Dataset\_Custom} 逻辑：训练集 70\%，验证集 10\%，测试集 20\%。

模型方面保留了三组：
\begin{itemize}
  \item 线性基线：\texttt{DLinear}
  \item 深度主模型：\texttt{PatchTST}
  \item 补充模型：\texttt{TimesNet}、\texttt{iTransformer}
\end{itemize}

其中，\texttt{PatchTST} 外汇核心组合使用 \texttt{10 epochs}，其余扩展实验使用已有结果目录继续做策略层回放，不重复训练。

\section{第二版策略设计}

\subsection{从单根开平到状态化持仓}

第一版回测是每根 K 线独立决策：预测下一根收益率超过阈值就开仓，下一根结束立即平仓。第二版改成状态化持仓：
\begin{itemize}
  \item 持仓后，如果同方向信号继续出现，则继续持有；
  \item 若信号转为空，不立即反向，允许继续持有或平仓，取决于参数设置；
  \item 若出现反向信号，则平旧仓并决定是否反手；
  \item 每次换手只在仓位变化时扣交易成本。
\end{itemize}

这一步直接对应反馈中的“如果预期连涨几根，中间就不要换手”。

\subsection{新增的策略参数}

第二版策略支持以下参数：
\begin{itemize}
  \item \texttt{threshold\_bps}：开仓阈值
  \item \texttt{hold\_until\_opposite}：是否持有到反向信号出现
  \item \texttt{min\_hold\_bars}：最短持有条数
  \item \texttt{stop\_loss\_bps}：止损阈值
  \item \texttt{cooldown\_bars}：平仓后冷却条数
  \item \texttt{ma\_window}：均线过滤窗口
  \item \texttt{long\_only}：仅做多
\end{itemize}

其中，均线过滤的实现方式是：价格位于均线之上才允许做多，价格位于均线之下才允许做空。这个过滤器已经实际跑过，但在当前样本上，它更像保守版本，而不是收益最高的主版本。

\section{结果对比}

\subsection{改进前后对照}

表 \ref{tab:before_after} 给出几组代表性组合的改进前后对照。改进前使用第一版口径；改进后使用第二版状态化持仓策略。

\begin{table}[htbp]
\centering
\caption{改进前后代表性结果对照}
\label{tab:before_after}
\small
\begin{tabular}{C{1.8cm} C{1.2cm} C{1.7cm} C{1.8cm} C{1.8cm} C{1.6cm}}
\toprule
品种 & 周期 & 模型 & 基线累计收益 & 改进后累计收益 & 收益改善 \\
\midrule
GBPUSD & M15 & PatchTST     & -12.01\% & +5.04\%  & +17.04\% \\
EURUSD & M15 & PatchTST     & -10.76\% & +4.36\%  & +15.12\% \\
GBPUSD & H1  & PatchTST     & -21.49\% & +4.01\%  & +25.49\% \\
EURUSD & H1  & PatchTST     & -20.59\% & +3.76\%  & +24.35\% \\
US500  & H1  & PatchTST     & -23.05\% & +22.05\% & +45.10\% \\
USTEC  & H4  & DLinear      & +7.49\%  & +35.40\% & +27.91\% \\
USTEC  & H4  & iTransformer & +7.54\%  & +37.74\% & +30.20\% \\
\bottomrule
\end{tabular}
\end{table}

改进后，收益改善主要来自两部分：一是阈值提高之后，交易数量明显下降；二是连续持仓把趋势段保留下来，减少了来回开平造成的摩擦。

\subsection{第二版策略下的代表性组合}

表 \ref{tab:improved} 展示第二版策略下最有代表性的正收益组合及其对应参数。

\begin{table}[htbp]
\centering
\caption{第二版策略下的代表性正收益组合}
\label{tab:improved}
\small
\begin{tabular}{C{1.7cm} C{1.1cm} C{1.7cm} C{1.1cm} C{1.0cm} C{1.0cm} C{1.6cm} C{1.5cm}}
\toprule
品种 & 周期 & 模型 & 阈值 & 交易数 & 止损 & 累计收益 & 最大回撤 \\
\midrule
GBPUSD & M15 & PatchTST     & 12bps & 10  & 40bps & +5.04\%  & 1.17\% \\
EURUSD & M15 & PatchTST     & 10bps & 8   & 无     & +4.36\%  & 2.33\% \\
GBPUSD & H1  & PatchTST     & 20bps & 8   & 40bps & +4.01\%  & 3.12\% \\
EURUSD & H1  & PatchTST     & 12bps & 29  & 20bps & +3.76\%  & 1.93\% \\
US500  & H1  & PatchTST     & 15bps & 145 & 40bps & +22.05\% & 3.97\% \\
USTEC  & H4  & DLinear      & 8bps  & 165 & 40bps & +35.40\% & 14.24\% \\
USTEC  & H4  & iTransformer & 20bps & 161 & 无     & +37.74\% & 7.19\% \\
\bottomrule
\end{tabular}
\end{table}

这些结果说明第二版策略已经不再是“高频开平导致大面积负收益”的状态，而是形成了可解释的正收益样本。

\subsection{各项修改的作用}

从参数搜索结果看，四项修改的作用并不完全相同：
\begin{enumerate}
  \item \textbf{提高阈值} 是最直接的降换手手段。外汇短周期从 \texttt{3bps} 提高到 \texttt{10--20bps} 后，交易次数大幅下降。
  \item \textbf{持有到反向信号} 是收益改善的核心驱动。多数正收益组合都使用了 \texttt{hold\_until\_opposite = True}。
  \item \textbf{止损} 对 \texttt{H1} 和指数组合帮助明显，能把回撤压下来。
  \item \textbf{均线过滤} 已经测试过。例如 \texttt{GBPUSD M15 + PatchTST + MA20} 可以得到约 \texttt{+0.82\%} 的收益、约 \texttt{0.33\%} 的最大回撤，但整体不如本轮最优参数组合强。
\end{enumerate}

\section{对题目的回答}

\subsection{哪个品种、哪个周期更适合深度学习}

如果“适合深度学习”定义为深度模型相对简单线性基线的提升空间，则第二版结论仍然主要落在外汇上：
\begin{center}
$\texttt{EURUSD} \approx \texttt{GBPUSD} > \texttt{USDJPY}$
\end{center}

周期排序为：
\begin{center}
\texttt{M15} \quad > \quad \texttt{H1} \quad > \quad \texttt{H4}
\end{center}

理由是外汇 \texttt{M15/H1} 仍然最能体现深度模型对线性基线的相对增益，且经过交易层改造后，已经能从负收益转成正收益。

\subsection{当前绝对收益最好的方向}

若按当前绝对收益衡量，则最强的样本已经不只限于外汇：
\begin{itemize}
  \item \texttt{USTEC H4 + iTransformer}：\(+37.74\%\)
  \item \texttt{USTEC H4 + DLinear}：\(+35.40\%\)
  \item \texttt{US500 H1 + PatchTST}：\(+22.05\%\)
\end{itemize}

因此，第二版的更准确表述是：
\begin{quote}
外汇 \texttt{M15/H1} 仍是最能体现深度模型价值的区域；指数低频场景在交易层修正之后则呈现出更强的绝对收益。
\end{quote}

\section{最终结论}

\begin{quote}
第二版实验将原先的单根 K 线开平策略改成了状态化持仓策略，并加入高阈值开仓、止损、冷却期和均线过滤。结果显示，这些改动足以把多组原本为负的组合拉成正收益。其中，\texttt{GBPUSD/EURUSD} 的 \texttt{M15/H1} 组合证明了深度模型在外汇短周期上的价值；\texttt{US500 H1} 与 \texttt{USTEC H4} 则说明交易层修正之后，指数场景也具备较强表现。就本题而言，提交物不再只是代码能力展示，而已经形成了可复现、可解释、并能够跑出正收益样本的研究结果。
\end{quote}

\vspace{0.8em}
\hrule
\vspace{0.7em}

{\small
说明：正文中引用的结果对应文件已落地在仓库 \texttt{quant/outputs} 与 \texttt{results} 目录中，可直接复核。}

\end{document}
